\clearpage
\VerdierPhysicalPageAnchor{0163}%
\begingroup\small\setlength{\parskip}{0pt}\setlength{\emergencystretch}{1.5em}
\VerdierSemanticAnchor{ch4-4-1-2-b}
b) 一个从 $\mathcal{F}\ell_2(\mathcal J)$ 到 $\mathcal D$ 中特异三角所成范畴的
函子，其形式为
\[
\omega_2=(\mathord\bullet\xrightarrow f\mathord\bullet\xrightarrow g\mathord\bullet)
\longmapsto
X_f\xrightarrow{X*u_1(\omega_2)}X_{gf}
\xrightarrow{X*u_2(\omega_2)}X_g
\xrightarrow{\delta(\omega_2)}X_f[1].
\]
为表示 $\mathcal J$ 型谱对象，我们将使用记号
\[
(X_f, f\in\operatorname{Ob}\mathcal{F}\ell_1(\mathcal J),\
\delta(\omega_2),
\ \omega_2\in\operatorname{Ob}\mathcal{F}\ell_2(\mathcal J)),
\]
在不会产生混淆时，也简记为 $(X_f,\delta)$。

\VerdierSemanticAnchor{ch4-4-1-3}
\medskip\noindent\textbf{4.1.3.}\quad
设 $C$ 为加法范畴，且
$F\colon\mathcal J\longrightarrow\operatorname{comp}(C)$ 为一个函子。
第 I 章 (3.3.6) 已说明如何为 $F$ 配置一个 $\mathcal J$ 型且取值于
$K(C)$ 的谱对象。第 I 章 (3.3.6) 所述的谱对象称为
\emph{与函子 $F$ 关联的谱对象}。

\VerdierSemanticAnchor{ch4-4-1-4}
\medskip\noindent\textbf{4.1.4.}\quad
设 $\mathcal A$ 为阿贝尔范畴，$\mathcal J$ 为一个范畴。一个
\emph{$\mathcal J$ 型且取值于 $\mathcal A$ 的谱对象}由以下资料组成：

a) 一列函子
$H^n\colon\mathcal{F}\ell_1(\mathcal J)\longrightarrow\mathcal A$，
$n\in\mathbb Z$。

b) 对每个对象
$\omega_2=(\mathord\bullet\xrightarrow f\mathord\bullet\xrightarrow g\mathord\bullet)$，
它属于 $\mathcal{F}\ell_2(\mathcal J)$，以及每个整数 $n$，给定一个态射
$\delta^n(\omega_2)\colon H^n(g)\to H^{n+1}(f)$，它关于 $\omega_2$
为函子性的，并使图式
\[
\cdots\longrightarrow H^n(f)
\xrightarrow{H^n*u_1(\omega_2)}H^n(gf)
\xrightarrow{H^n*u_2(\omega_2)}H^n(g)
\xrightarrow{\delta^n(\omega_2)}H^{n+1}(f)
\longrightarrow\cdots
\]
成为 $\mathcal A$ 中的无界正合列。

我们将使用记号
\[
(H^n(f), f\in\operatorname{Ob}\mathcal{F}\ell_1(\mathcal J),\
\delta^n(\omega_2),
\ \omega_2\in\operatorname{Ob}\mathcal{F}\ell_2(\mathcal J))
\]
表示 $\mathcal J$ 型谱对象，在不会产生混淆时，也简记为
$(H^n(f),\delta)$。

\VerdierSemanticAnchor{ch4-4-1-5}
\medskip\noindent\textbf{4.1.5.}\quad
设 $F\colon\mathcal D\to\mathcal D'$ 为两个三角范畴之间的正合函子，
$(X_f,\delta)$ 为一个 $\mathcal J$ 型且取值于 $\mathcal D$ 的谱对象。
函子对 $(FX_f,F\delta)$ 是一个 $\mathcal J$ 型且取值于 $\mathcal D'$ 的
谱对象。类似地，从三角范畴到阿贝尔范畴的上同调函子
$H\colon\mathcal D\to\mathcal A$，把取值于 $\mathcal D$ 的谱对象变为
取值于 $\mathcal A$ 的谱对象。

\VerdierSemanticAnchor{ch4-4-1-6}
\medskip\noindent\textbf{4.1.6.}\quad
设 $\mathcal D$ 为三角范畴，$(X_f,\delta)$、$(Y_f,\delta')$ 为两个
$\mathcal J$ 型且取值于 $\mathcal D$ 的谱对象。一个\emph{谱对象的态射}
是一个函子态射
\[
m_f\colon X_f\longrightarrow Y_f,\qquad
f\in\operatorname{Ob}\mathcal{F}\ell_1(\mathcal J),
\]
\endgroup
